% BHU PYQ -- Manually transcribed & error-corrected
% Paper: BCF-116 : Business Statistics
% Exam: B.Com. (Hons.) FMM Semester I Examination, 2016-17

\documentclass[11pt,a4paper]{article}
\usepackage[a4paper,top=2.5cm,bottom=2.5cm,left=2.8cm,right=2.8cm,headheight=14pt]{geometry}
\usepackage{amsmath,amssymb}
\usepackage[T1]{fontenc}
\usepackage[utf8]{inputenc}
\usepackage{lmodern,microtype,fancyhdr,enumitem,hyperref,lastpage,parskip}

\hypersetup{
    colorlinks=true,
    urlcolor=black,
    linkcolor=black,
    pdftitle={BCF-116: Business Statistics -- BHU B.Com. (Hons.) FMM Sem I 2016-17}
}

\pagestyle{fancy}
\fancyhf{}
\renewcommand{\headrulewidth}{0.4pt}
\renewcommand{\footrulewidth}{0.4pt}
\fancyhead[L]{\small   \textit{Banaras Hindu University}}
\fancyhead[R]{\small   \textit{BCF-116: Business Statistics}}
\fancyfoot[C]{\small Page  \thepage\ of \pageref{LastPage}}


\newlist{parts}{enumerate}{2}
\setlist[parts,1]{label=(\alph*),leftmargin=*,itemsep=5pt,topsep=3pt}
\setlist[parts,2]{label=(\roman*),leftmargin=*,itemsep=3pt,topsep=2pt}

\newcommand{\pts}[1]{\hfill   \textit{[#1]}}


\begin{document}


\begin{center}
  {\large\bfseries BANARAS HINDU UNIVERSITY}\\[2pt]
  {
\normalsize B.Com. (Hons.) FMM --- Semester I Examination, 2016-17}\\[6pt]
  \rule{0.8\linewidth}{0.5pt}\\[6pt]
  {\large\bfseries Paper No. BCF-116: Business Statistics}\\[4pt]
  {
\normalsize Time: 3 Hours\hfill Maximum Marks: 70}
\end{center}

\rule{\linewidth}{0.4pt}

\smallskip



\noindent   \textit{Instructions: Answer all questions. All questions carry equal marks (14 marks each).}

\bigskip




\noindent   \textbf{Question 1.}
What is statistics? Explain the limitations and importance of statistics in the present time.\pts{14}

\centerline{   \textbf{OR}}

What is primary data? What methods are used to collect it? Explain their relative merits and demerits.\pts{14}

\medskip\rule{\linewidth}{0.2pt}




\noindent   \textbf{Question 2.}
An aeroplane covered a distance of 800 kilometers with four different speeds of 100, 200, 300 and 400 km. per hour for the first, second, third and fourth quarters of the distance respectively. Find the average speed in km. per hour and test the validity of your answer.\pts{14}

\centerline{   \textbf{OR}}

Find out the mode and median wages from the following data:\pts{14}

\begin{center}
  \renewcommand{\arraystretch}{1.2}
  
\begin{tabular}{|l|c|c|c|c|c|}
    \hline
   \textbf{Wages (Rs.)} & 0--10 & 10--20 & 20--30 & 30--40 & 40--50 \\ \hline
   \textbf{No. of Workers} & 5 & 7 & 15 & 25 & 8 \\ \hline
  \end{tabular}
\end{center}

\medskip\rule{\linewidth}{0.2pt}




\noindent   \textbf{Question 3.}
Goals scored by teams A and B in a football season were as follows:

\begin{center}
  \renewcommand{\arraystretch}{1.2}
  
\begin{tabular}{|c|c|c|}
    \hline
   \textbf{No. of Goals (Scored in a match)} & \multicolumn{2}{c|}{   \textbf{No. of Matches}} \\ \cline{2-3}
    &   \textbf{Team A} &   \textbf{Team B} \\ \hline
    0 & 27 & 17 \\ \hline
    1 & 9 & 9 \\ \hline
    2 & 8 & 6 \\ \hline
    3 & 5 & 5 \\ \hline
    4 & 4 & 3 \\ \hline
  \end{tabular}
\end{center}
By calculating the Coefficient of Variation for each, find which team may be considered more consistent.\pts{14}

\centerline{   \textbf{OR}}

If the values of Mean, Median and Mode are Rs. 33.26, Rs. 33.5 and Rs. 34 respectively in the following distribution, calculate the missing frequencies:\pts{14}

\begin{center}
  \renewcommand{\arraystretch}{1.2}
  
\begin{tabular}{|l|c|c|c|c|c|c|c|c|}
    \hline
   \textbf{Wages (Rs.)} & 0--10 & 10--20 & 20--30 & 30--40 & 40--50 & 50--60 & 60--70 &   \textbf{Total} \\ \hline
   \textbf{Frequency} & 8 & 32 & ? & ? & ? & 12 & 8 & 460 \\ \hline
  \end{tabular}
\end{center}

\medskip\rule{\linewidth}{0.2pt}




\noindent   \textbf{Question 4.}
Explain clearly the meaning of ``Time Series Analysis''. Indicate the importance of such analysis in business.\pts{14}

\centerline{   \textbf{OR}}

Below are given the figures of milk production (in thousand litres) of a milk dairy:

\begin{center}
  \renewcommand{\arraystretch}{1.2}
  
\begin{tabular}{|l|c|c|c|c|c|c|c|}
    \hline
   \textbf{Year} & 2008 & 2009 & 2010 & 2011 & 2012 & 2013 & 2014 \\ \hline
   \textbf{Production} & 77 & 88 & 94 & 85 & 91 & 98 & 90 \\ \hline
  \end{tabular}
\end{center}

\begin{parts}
  \item Fit a straight line trend by the least squares method.
  \item What is the monthly increase in the production of the milk dairy?
  \item Estimate the production for the year 2018.
\end{parts}\pts{14}

\medskip\rule{\linewidth}{0.2pt}




\noindent   \textbf{Question 5.}
What is an index number? Examine the various problems involved in the construction of index numbers. Also discuss the uses of index numbers.\pts{14}

\centerline{   \textbf{OR}}

The following table gives the average wholesale prices of three groups of commodities for the years 2011 to 2015. Compute the chain base index numbers based on 2011:\pts{14}

\begin{center}
  \renewcommand{\arraystretch}{1.2}
  
\begin{tabular}{|c|c|c|c|c|c|}
    \hline
   \textbf{Group} &   \textbf{2011} &   \textbf{2012} &   \textbf{2013} &   \textbf{2014} &   \textbf{2015} \\ \hline
    I   & 2 & 3 & 5 & 7 & 8 \\ \hline
    II  & 8 & 10 & 12 & 14 & 18 \\ \hline
    III & 4 & 5 & 7 & 9 & 12 \\ \hline
  \end{tabular}
\end{center}

\vfill
\rule{\linewidth}{0.4pt}

\begin{center}\small
\href{https://bhu-exam.vercel.app/}{Click here to visit our website}\\[4pt]
\href{https://me-aryan.vercel.app/}{Click here to visit our contributor}\\[4pt]
\href{https://quashan-me.vercel.app/}{Click here to visit our contributor}
\end{center}

\end{document}
